%! TEX program = xelatex
% ============================================================
% 硕士学位论文主文档
% 编译顺序：xelatex -> biber -> xelatex -> xelatex
% 也可直接：latexmk（已提供 .latexmkrc）
% ============================================================
\documentclass[
  degree   = master,   % 硕士论文（bachelor / master / doctor）
  zihao    = -4,       % 小四号字（12pt，正文基准字号）
  fontset  = auto,     % 按操作系统自动选择中文字体（win/mac/linux）
  openright,           % 章节从奇数页开始
  UTF8
]{ctexbook}

% 全局宏包与设置
% ============================================================
% setup/preamble.tex —— 宏包与全局设置
% 本文件由 main.tex 通过 \input 引入，不要单独编译。
% 原则：只放"依赖与全局开关"，不放具体格式与内容。
% ============================================================

% ---------- 数学与符号 ----------
\usepackage{amsmath, amssymb, amsthm}
\usepackage{mathtools}
\usepackage{bm}                 % 粗体数学符号

% ---------- 图表 ----------
\usepackage{graphicx}           % 插图
\graphicspath{{figures/}{tables/}}  % 图片/表格搜索路径
\usepackage{subcaption}         % 子图（subfigure）
\usepackage{booktabs}           % 三线表
\usepackage{multirow}           % 表格跨行
\usepackage{makecell}           % 单元格内换行
\usepackage{longtable}          % 跨页长表
\usepackage{diagbox}            % 斜线表头（按需）
\usepackage{array}

% ---------- 页面尺寸与页眉页脚 ----------
\usepackage{geometry}           % 纸张与页边距
\usepackage{fancyhdr}           % 页眉/页脚定制（中国地质大学格式）

% ---------- 算法与代码 ----------
\usepackage{algorithm}
\usepackage{algorithmicx}
\usepackage{algpseudocode}
\usepackage{listings}           % 代码 listings
\usepackage{xcolor}

\lstset{
  basicstyle   = \small\ttfamily,
  breaklines   = true,
  frame        = single,
  numbers      = left,
  numberstyle  = \tiny,
  keywordstyle = \color{blue!70!black},
  commentstyle = \color{green!50!black},
  stringstyle  = \color{orange!80!black},
  captionpos   = b
}

% ---------- 参考文献（GB/T 7714-2005 国标，见规范 2.6）----------
% 顺序编码制：序号以方括号标注，置于引用处右上角（上标）。
% 正文引用请用 \supercite{key}（上标）；gb7714-2005 同时支持 \cite（行内）。
\usepackage[
  backend  = biber,
  style   = gb7714-2005,
  gbalign = left,
  sorting = none
]{biblatex}
\addbibresource{bib/references.bib}

% ---------- 交叉引用与超链接 ----------
\usepackage{hyperref}           % 建议最后加载此类宏包
\hypersetup{
  colorlinks = true,
  linkcolor  = black,
  citecolor  = black,
  urlcolor   = black,
  pdfborder  = {0 0 0}
}
\usepackage{cleveref}           % 智能引用（在 hyperref 之后）
\usepackage{siunitx}            % 单位与数字排版（规范 2.7：GB 3100~3102）
  \sisetup{
    inter-unit-product = \ensuremath{{}\cdot{}},  % 单位间用点乘
    per-mode = symbol,                             % 用 / 表示"每"
    separate-uncertainty = true
  }
\usepackage{enumitem}           % 列表环境定制

% 字体设置（英文字体：Times New Roman）
% ============================================================
% setup/fonts.tex —— 字体设置
%
% 要求：英文字体 = Times New Roman（新罗马）
% 中文：沿用 ctex 根据操作系统自动选择的字体
%       （Windows: 宋体/黑体/楷体/仿宋；macOS: 华文系列；Linux: Fandol）
%
% 说明：\setmainfont 只改变"西文（英文）衬线字体"，
%       中文字体仍由 ctex 的 \setCJKmainfont 控制，两者互不干扰，
%       因此英文显示为新罗马、中文保持规范字体。
% ============================================================

% —— 西文字体 ——
\setmainfont{Times New Roman}     % 英文正文衬线 = Times New Roman
\setsansfont{Arial}               % 西文无衬线（标题/图表标签可选）
\setmonofont{Courier New}         % 西文等宽（代码、算法）

% —— 辅助命令：在中文（宋体/黑体）环境中强制西文/数字用 Times New Roman ——
% 例：\tnr{2020110101} 用于封面"学号"等需要 TNR 阿拉伯数字处。
\newcommand{\tnr}[1]{{\fontspec{Times New Roman}#1}}

% —— 如需手动指定中文字体（例如 Linux 无系统字体时）——
% 取消下面注释并改为本机可用字体名，同时把 main.tex 的
% fontset=auto 改为 fontset=none：
% \setCJKmainfont{SimSun}          % 宋体（正文）
% \setCJKsansfont{SimHei}          % 黑体（无衬线/部分标题）
% \setCJKfamilyfont{zhkai}{KaiTi}  % 楷体
% \setCJKfamilyfont{zhfs}{FangSong} % 仿宋

% 自定义命令、定理环境、算法汉化
% ============================================================
% setup/commands.tex —— 自定义命令、缩写、定理环境、算法汉化
% ============================================================

% ---------- 常用英文缩写 ----------
\newcommand{\ie}{\textit{i.e.}}
\newcommand{\eg}{\textit{e.g.}}
\newcommand{\etal}{\textit{et al.}}

% ---------- 向量/矩阵粗体 ----------
\newcommand{\vect}[1]{\bm{#1}}
\newcommand{\mat}[1]{\bm{#1}}

% ---------- 数学算子 ----------
\DeclareMathOperator*{\argmin}{arg\,min}
\DeclareMathOperator*{\argmax}{arg\,max}

% ---------- 定理/定义/引理环境 ----------
\newtheorem{theorem}{定理}[chapter]
\newtheorem{lemma}{引理}[chapter]
\newtheorem{definition}{定义}[chapter]
\newtheorem{proposition}{命题}[chapter]
\newtheorem{remark}{注}[chapter]

% ---------- 算法环境关键词汉化（algpseudocode）----------
\renewcommand{\algorithmicrequire}{\textbf{输入}}
\renewcommand{\algorithmicensure}{\textbf{输出}}
\renewcommand{\algorithmicfor}{\textbf{对于}}
\renewcommand{\algorithmicdo}{\textbf{做}}
\renewcommand{\algorithmicwhile}{\textbf{当}}
\renewcommand{\algorithmicthen}{\textbf{则}}
\renewcommand{\algorithmicif}{\textbf{如果}}
\renewcommand{\algorithmicelse}{\textbf{否则}}
\renewcommand{\algorithmicreturn}{\textbf{返回}}

% 正文排版规范（标题/段落/图表题/公式编号，按学位论文要求）
% ============================================================
% setup/format.tex —— 正文排版规范（按学位论文要求集中配置）
%
% 本文件由 main.tex 通过 \input 引入（放在 commands 之后）。
% 覆盖内容：
%   1) 章节/节/小节标题：字号、字体、对齐、上下间距、序号与题名间距
%   2) 正文段落：宋体小四、英文 Times New Roman、两端对齐、
%                首行缩进 2 汉字符、固定行距 20 磅、段前段后 0
%   3) 图/表序号与题名：宋体五号、居中、单倍行距、序号与题名空 1 汉字符
%   4) 表达式编号：右对齐、(n) 形式、Times New Roman 五号
%
% 字号对照：三号≈16pt / 四号≈14pt / 小四≈12pt / 五号≈10.5pt
% ============================================================

% ---------- 1. 标题格式（ctex 标题机制）----------
\ctexset{
   % 各章标题：黑体三号(16pt)（黑体本身即粗，不再叠加 \bfseries 以免伪粗警告）居中，单倍行距，上下空 2 行
  % 章序号用中文数字（第一章），序号与题名之间空一个汉字符
  chapter = {
    name      = {第,章},
    number    = \chinese{chapter},   % 若学校要求"第1章"阿拉伯数字，改为 \arabic{chapter}
    format    = \centering\heiti\zihao{3},
    aftername = \quad,               % 序号与题名间空一个汉字符
    beforeskip = 2\baselineskip,     % 上空约 2 行
    afterskip  = 2\baselineskip,     % 下空约 2 行
    fixskip   = true,                % 固定间距（无弹性）
  },
  % 一级节标题：黑体四号(14pt)（黑体不加 \bfseries）居中，单倍行距，上下空 1 行
  section = {
    format    = \centering\heiti\zihao{4},
    aftername = \quad,
    beforeskip = \baselineskip,      % 上空约 1 行
    afterskip  = \baselineskip,      % 下空约 1 行
    fixskip   = true,
  },
  % 二级节标题：黑体小四(12pt)居左，单倍行距，上下空 1 行
  subsection = {
    format    = \raggedright\heiti\zihao{-4},
    aftername = \quad,
    beforeskip = \baselineskip,
    afterskip  = \baselineskip,
    fixskip   = true,
  },
  % 三级节标题（如使用）：规范表3.8未单独规定，按 CY/T 35-2001 顺延为
  % 黑体小四、顶格（序号与题名间不额外加空格），上下空约半行
  subsubsection = {
    format    = \raggedright\heiti\zihao{-4},
    aftername = {},
    beforeskip = 0.5\baselineskip,
    afterskip  = 0.5\baselineskip,
    fixskip   = true,
  },
}

% ---------- 2. 正文段落 ----------
% 宋体小四由文档类 zihao=-4 统一控制；英文 Times New Roman 由 fonts.tex 控制
% 其余：两端对齐、首行缩进 2 汉字符、固定行距 20 磅、段前段后 0
\setlength{\parindent}{2\ccwd}      % 首行左缩进 2 个汉字符
\setlength{\parskip}{0pt}           % 段前段后 0（段与段之间不额外留空）
% 固定值行距 20 磅：直接控制 \baselineskip。
% 注意 1：setspace 并没有 \exactlineskipmode 命令，也没有 [exact] 选项
%         （会报 "Unknown option `exact'"），不可依赖它。
% 注意 2：直接改 \baselineskip 会被字号命令（\zihao{5} 等）重置。
% 方案：用 ctex 官方钩子 \ctex_patch_cmd 把"固定 20pt"追加进各字号的
%       内部定义——每次切换字号都会重新执行它们，行距恒定。
%       \lineskiplimit=-1000pt 是 TeX"固定行距"模式，等效 Word"固定值"。
\newcommand{\fixLineskip}[1]{%
  \setlength{\baselineskip}{#1}%
  \setlength{\lineskip}{\z@}%
  \setlength{\lineskiplimit}{-1000\p@}%  % TeX 固定行距模式
}
\makeatletter
% \normalsize/\small/\footnotesize 分别调用内部 \@normalsize/\@small/\@footnotesize。
% LaTeX 内核的 \addto 要求第二参数是"已展开的命令名"，因此先把
% \fixLineskip{20pt} 展开到 \next（用 \dimexpr...\relax 隔离 \p@ 的 @），再追加。
\expandafter\def\expandafter\next\expandafter{\expandafter\fixLineskip\expandafter{\dimexpr20\p@\relax}}
\addto\@normalsize{\next}
\@ifundefined{@small}{}{\addto\@small{\next}}
\@ifundefined{@footnotesize}{}{\addto\@footnotesize{\next}}
\makeatother
\AtBeginDocument{\normalsize}   % 使上述设置立即生效

% ---------- 3. 图 / 表序号与题名 ----------
% 宋体五号(10.5pt)居中，单倍行距，序号与题名之间空一个汉字符
% 图序置于图下方（默认），表序置于表上方
\usepackage{caption}
% 注意：caption 的 font 键不接受 \songti、\zihao{5} 这类"声明式"命令，
% 必须使用 caption 提供的受控字体/字号名（标准尺寸名或 \DeclareCaptionFont
% 自定义注册名）。中文规范"题注宋体五号居中"实现如下：
%   wuhao  = 严格五号（10.5pt/13pt 单倍行距）；
%   zhsongcn = ctex 定义的宋体族控制序列 \songti。caption 的 font 键
%            按名字展开字体命令，因此用 \csname songti\endcsname 间接
%            引用；songti 由 ctex（scheme=chinese）定义。
\DeclareCaptionFont{wuhao}{\fontsize{10.5pt}{13pt}\selectfont}
% 宋体族名：统一使用 ctex 定义的 CJKfamily "zhsong"
% ctex 定义的 CJKfamily "zhsong"（windows/mac/linux/fandol 均注册此名）。
% caption 在正常文档阶段展开，此时 \songti 已由 ctex 定义。
\DeclareCaptionFont{zhsongcn}{\csname songti\endcsname}  % 宋体（ctex 提供）
\captionsetup{
  font         = {zhsongcn,wuhao},      % 宋体五号
  labelfont    = {zhsongcn,wuhao},      % 图/表序号同样宋体五号
  justification = centering,          % 居中
  labelsep     = quad,                % 序号与题名之间空一个汉字符
  skip         = 6pt,                 % 图/表与题名的间距
}
\captionsetup[table]{position=above}   % 表序在表上方
\captionsetup[figure]{position=below}  % 图序在图下方

% ---------- 4. 表达式编号 ----------
% 右对齐、(n) 形式、Times New Roman 五号（默认已带圆括号，仅改字体与字号）
\makeatletter
\def\maketag@@@#1{\hbox{\m@th\normalfont\zihao{5}#1}}
\makeatother

% ---------- 5. 目录排版 ----------
% 标题"目录"：沿用 ctexset 的 chapter 样式（黑体三号加粗居中，单倍行距），无需额外设置。
% 目录条目要求：
%   各章目录  ：宋体四号(14pt)，固定行距 20 磅，两端对齐，页码右对齐，左缩进 0
%   一级节    ：宋体小四(12pt)，固定行距 20 磅，两端对齐，页码右对齐，左缩进 1 汉字符
%   二级节    ：宋体小四(12pt)，固定行距 20 磅，两端对齐，页码右对齐，左缩进 2 汉字符
\usepackage{titletoc}
\titlecontents{chapter}[0pt]{%                        % 左缩进 0
    \zihao{4}\songti\baselineskip=20pt%               % 宋体四号 + 固定 20 磅行距
  }{%
    \thecontentslabel\quad%                           % 编号（第一章）+ 空一个汉字符
  }{%                                                  % 无编号条目（如参考文献/致谢）
    {}%
  }{%
    \hfill\contentspage%                               % 页码右对齐（无点线）
  }
\titlecontents{section}[1\ccwd]{%                      % 左缩进 1 汉字符
    \zihao{-4}\songti\baselineskip=20pt%              % 宋体小四 + 固定 20 磅行距
  }{%
    \thecontentslabel\quad%
  }{%
    {}%
  }{%
    \hfill\contentspage%
  }
\titlecontents{subsection}[2\ccwd]{%                   % 左缩进 2 汉字符
    \zihao{-4}\songti\baselineskip=20pt%
  }{%
    \thecontentslabel\quad%
  }{%
    {}%
  }{%
    \hfill\contentspage%
  }

% 图目录 / 表目录：标题沿用 chapter 样式（黑体三号加粗居中）；
% 条目按一级节样式：宋体小四(12pt)、固定 20 磅、页码右对齐、左缩进 1 汉字符
\titlecontents{figure}[1\ccwd]{%
    \zihao{-4}\songti\baselineskip=20pt%
  }{%
    \thecontentslabel\quad%
  }{}{%
    \hfill\contentspage%
  }
\titlecontents{table}[1\ccwd]{%
    \zihao{-4}\songti\baselineskip=20pt%
  }{%
    \thecontentslabel\quad%
  }{}{%
    \hfill\contentspage%
  }

% ---------- 6. 中英文摘要 ----------
% 标题：中文"摘要"黑体小二(18pt)加粗居中；英文"Abstract" Times New Roman 小二(18pt)加粗居中；单倍行距
% 正文：中文宋体小四 / 英文 TNR 小四，固定 20 磅，段前段后 0（全局已设）
% 关键词："关键词"三字加粗（黑体）/ "Key Words"两词加粗（TNR），字号同正文
\newcommand{\cnabstracttitle}{%
  \begin{center}\heiti\zihao{-2} 摘要\end{center}%
  \vspace{0.5\baselineskip}\zihao{-4}\songti\noindent\ignorespaces
}
\newcommand{\enabstracttitle}{%
  \begin{center}\bfseries\zihao{-2} Abstract\end{center}%
  \vspace{0.5\baselineskip}\normalfont\zihao{-4}\noindent\ignorespaces
}

% ---------- 7. 纸张与页眉页脚（中国地质大学硕士论文格式）----------
% 纸张：A4（210×297），白色
% 页边距：上/下 3cm，左/右 3cm，装订线 0cm
% 页眉距顶 2.5cm（headheight=2.5cm，headsep=0.5cm）；页脚距底 2.0cm（footskip=2.0cm）
\geometry{
  a4paper,
  top        = 3cm,
  bottom     = 3cm,
  left       = 3cm,
  right      = 3cm,
  headheight = 2.5cm,   % 页眉区域高度（页眉底边距顶 2.5cm，贴顶）
  headsep    = 0.5cm,   % 页眉底边到正文的距离
  % 注意：geometry 没有 footheight 键（正确键名为 bottomheight / footskip），
  % 原 footheight=1.0cm 属无效选项会被忽略并警告，此处删除；
  % 页脚位置由 footskip 控制。
  footskip   = 2.0cm,   % 正文底边到页脚基线的距离（页脚距底约 2.0cm，贴底）
  bindingoffset = 0cm   % 装订线 0cm
}

% 页眉页脚样式：正文（第一章起）使用 running header
%   奇数页居中：中国地质大学<学位类别>
%   偶数页居中：作者姓名：论文题目
%   页码外侧（奇→右、偶→左），宋体五号(10.5pt)
%   页眉之下有下划线（headrule）
\fancyhf{}                                     % 清空默认页眉页脚
\renewcommand{\headrulewidth}{0.4pt}           % 页眉下划线
\renewcommand{\footrulewidth}{0pt}             % 页脚无线
\fancyhead[CO]{\songti\zihao{5} 中国地质大学\thesisCategory}      % 奇数页：学位类别
\fancyhead[CE]{\songti\zihao{5} \thesisAuthor：\thesisTitle}      % 偶数页：作者：题目
\fancyfoot[RO]{\songti\zihao{5} \thepage}      % 奇数页页码在右侧（外侧）
\fancyfoot[LE]{\songti\zihao{5} \thepage}      % 偶数页页码在左侧（外侧）

% 英文摘要页眉：按规范"Abstract 部分用 Times New Roman 五号"，标题居中（TNR）
\fancypagestyle{abstracten}{%
  \fancyhf{}%
  \renewcommand{\headrulewidth}{0.4pt}%
  \fancyhead[C]{\zihao{5} \thesisTitleEN}%      % 英文标题（TNR，主字体即 Times New Roman）
  \fancyfoot[C]{\songti\zihao{5} \thepage}%     % 页码仍按规范用宋体五号
}


% ============================================================
% 论文基本信息（按需修改）
% ============================================================
\newcommand{\thesisTitle}{多无人机智能调度中的任务分配与航迹规划方法研究}
\newcommand{\thesisTitleEN}{Task Allocation and Path Planning Methods for Multi-UAV Intelligent Scheduling}
\newcommand{\thesisAuthor}{张佳齐}
\newcommand{\thesisAuthorEn}{Jiaqi Zhang}  % 英文姓名用拼音（姓在前可改为 Zhang Jiaqi，按学校要求）
\newcommand{\thesisAdvisor}{郑坤\ 教授}
\newcommand{\thesisSchool}{中国地质大学（武汉）}
\newcommand{\thesisCategory}{硕士学位论文}% 学位类别（页眉奇数页用，学术型硕士=硕士学位论文）：博士学位论文 / 同等学力博士学位论文 / 硕士学位论文 / 硕士专业学位论文(全日制) / 硕士专业学位论文(非全日制) / 同等学力硕士学位论文
\newcommand{\thesisMajor}{地理学}              % 一级学科（按实际专业修改）
\newcommand{\thesisResearchDir}{（按实际研究方向填写）}% 研究方向
\newcommand{\thesisOrg}{地理与信息工程学院}      % 培养单位（所在学院全称）
\newcommand{\thesisSchoolCode}{10491}          % 学校代码（教育部批准，中国地质大学(武汉)）
\newcommand{\thesisStudentID}{1202510805}      % 学号
\newcommand{\thesisCLC}{（填写）}              % 中图分类号（封面用，按《中国图书馆分类法》第4版）
\newcommand{\thesisUDC}{（填写）}              % UDC（按《国际十进分类法》）
\newcommand{\thesisSecretLevel}{公开}          % 密级（公开论文可标可不标）
\newcommand{\thesisDate}{二〇二六年六月}        % 完成时间（汉字，不用阿拉伯数字，见规范 3.2/3.4）

\begin{document}

% ============================================================
% 前置部分（按学位论文组成顺序：封面→题名页→声明/承诺书/授权书→作者简介→答辩名单→中英文摘要→目录→图表清单）
% ============================================================
\frontmatter
\pagenumbering{roman}          % 前置部分用罗马数字（第一章起改为阿拉伯数字）
\pagestyle{plain}              % 前置部分：页脚居中罗马页码、无页眉
% ============================================================
% frontmatter/cover.tex —— 中文封面（规范 3.2）
% 字体字号严格按《中国地质大学（武汉）研究生学位论文写作规范》：
%   学校名 / 学位类型：宋体 26 磅(一号) 加粗
%   论文题目：黑体 22 磅(二号) 加粗居中，单倍行距（可分两行）
%   学号：Times New Roman 16 磅(三号) 加粗
%   论文作者/学科专业/指导教师/培养单位：宋体 16 磅(三号) 加粗
%   日期：宋体 16 磅(三号) 汉字居中，不用阿拉伯数字
%   分类号/密级/UDC/学校代码：封面著录项（见规范 1.1）
% ============================================================
\begin{titlepage}
  \centering

  % ---- 著录信息（封面右上角小字）----
  \vspace*{-1.2cm}
  \begin{flushleft}
    \zihao{5}\songti
    分类号：\underline{\thesisCLC}\hfill 学校代码：\tnr{\thesisSchoolCode}\\[6pt]
    密\hspace{1em}级：\underline{\thesisSecretLevel}\hfill UDC：\underline{\thesisUDC}
  \end{flushleft}

  \vspace{1.6cm}
  {\songti\zihao{1}\textbf{\thesisSchool}}\\[0.4cm]
  {\songti\zihao{1}\textbf{\thesisCategory}}\\[1.8cm]

  {\heiti\zihao{2}\textbf{\thesisTitle}}\\[2cm]

  \begin{tabular}{rl}
    \zihao{3}\songti\bfseries 学\hspace{1em}号： & \zihao{3}\tnr{\textbf{\thesisStudentID}} \\[0.5cm]
    \zihao{3}\songti\bfseries 论文作者： & \zihao{3}\songti\bfseries \thesisAuthor \\[0.5cm]
    \zihao{3}\songti\bfseries 学科专业： & \zihao{3}\songti\bfseries \thesisMajor \\[0.5cm]
    \zihao{3}\songti\bfseries 指导教师： & \zihao{3}\songti\bfseries \thesisAdvisor \\[0.5cm]
    \zihao{3}\songti\bfseries 培养单位： & \zihao{3}\songti\bfseries \thesisOrg \\[0.5cm]
  \end{tabular}

  \vfill
  {\zihao{3}\songti \thesisDate}
\end{titlepage}
              % 1. 中文封面（titlepage 环境自动无页码）
% ============================================================
% frontmatter/titlepage_cn.tex —— 中文题名页（规范 3.4）
% 【V4 校准】严格按“附件1-11 p.1（学术型硕士式样）”重建：
%   顺序（顶→底）：
%     学校代码/研究生学号 → 中国地质大学 → 硕士学位论文
%     → 论文题目 → 姓名 → 学科专业 → 指导教师 → 培养单位 → 完成日期
%   此前（V3）误用了师兄“硕士专业学位”题名页的顺序，与学术型附件相反，
%   本次已纠正为附件学术型顺序。
%   字体：校名/类别 宋体一号加粗；论文题目 黑体二号加粗；信息块 宋体三号；
%        学号等阿拉伯数字用 Times New Roman（\tnr）。
%   标签经 \labelfour 等宽分散（与附件式样一致：2 字大间距、4 字小间距）。
%   注：沿用 tabular（fixLineskip 修复后正常分行、顺序正确），不塌行。
% ============================================================
\clearpage
\thispagestyle{empty}

\begin{center}
  % —— 整体下移，使顶部“学校代码”落在距页顶约 5.8cm（对齐附件学术型式样）——
  \vspace*{2.8cm}
  % —— 顶部：学校代码 / 研究生学号（左锚 r0≈0.019）——
  \zihao{3}\songti
  \makebox[\textwidth][l]{%
    \hspace*{0.019\textwidth}%
    学校代码：\tnr{\thesisSchoolCode}\quad 研究生学号：\tnr{\thesisStudentID}}

  \vspace*{1.1cm}
  % —— 学位授予单位名称（宋体一号加粗居中）——
  {\songti\zihao{1}\textbf{\thesisSchool}}

  \vspace*{1.1cm}
  % —— 学位论文类别（宋体一号加粗居中）——
  {\songti\zihao{1}\textbf{\thesisCategory}}

  \vspace*{1.3cm}
  % —— 论文题目（黑体二号加粗居中，超宽自动折行；块宽约 0.70 版心使两行均衡）——
  \parbox{0.70\textwidth}{\centering\heiti\zihao{2}\thesisTitle}

  \vspace*{1.7cm}
  % —— 信息块（左锚 r0≈0.268；顺序按附件：姓名→学科专业→指导教师→培养单位）——
  \makebox[\textwidth][l]{%
    \hspace*{0.268\textwidth}%
    \begin{tabular}{rl}
      \labelfour{姓名}：     & \zhlabelsep\thesisAuthor   \\[0.42cm]
      \labelfour{学科专业}： & \zhlabelsep\thesisMajor   \\[0.42cm]
      \labelfour{指导教师}： & \zhlabelsep\thesisAdvisor \\[0.42cm]
      \labelfour{培养单位}： & \zhlabelsep\thesisOrg
    \end{tabular}}

  \vspace*{3.2cm}
  % —— 完成日期（汉字，居中）——
  \thesisDate
\end{center}
       % 2. 中文题名页
% ============================================================
% frontmatter/titlepage_en.tex —— 英文题名页（规范 3.5）
% 英文排版要求（全文 Times New Roman）：
%   A Dissertation Submitted to China University of Geosciences：16 磅(三号)
%   Title（论文题目英文）：22 磅(二号) 加粗居中，单倍行距
%   Master Candidate / Major / Supervisor：16 磅(三号) 居中
%   China University of Geosciences, Wuhan 430074, P.R. China：16 磅(三号) 居中
% 注：英文题名页不另写学位类型，由 "A Dissertation Submitted ..." 体现。
% ============================================================
\clearpage
\thispagestyle{empty}

\begin{center}
  \zihao{3}
  A Dissertation Submitted to China University of Geosciences

  \vspace{1.2cm}
  {\zihao{2}\bfseries \thesisTitleEN}

  \vspace{1.4cm}
  Master Candidate: \thesisAuthorEn \\[0.5cm]
  Major: Geography \\[0.5cm]
  Supervisor: \thesisAdvisorEn

  \vspace{1.6cm}
  China University of Geosciences, Wuhan 430074, P.R. China
\end{center}
       % 3. 英文题名页
% ============================================================
% frontmatter/declaration.tex —— 学位论文原创性声明（第 4 项）
% 通用模板，请以学校下发的最新版本为准替换正文。
% 注：导师承诺书见 supervisor_commitment.tex，使用授权书见 authorization.tex
% ============================================================
\chapter*{学位论文原创性声明}
\addcontentsline{toc}{chapter}{学位论文原创性声明}

本人郑重声明：所呈交的学位论文，是本人在导师指导下，独立进行研究工作所取得的成果。
除文中已经注明引用的内容外，本论文不包含任何其他个人或集体已经发表或撰写过的作品成果。
对本文的研究作出重要贡献的个人和集体，均已在文中以明确方式标明。
本人完全意识到本声明的法律结果由本人承担。

\vspace{1.5cm}
\noindent 学位论文作者签名：\underline{\hspace{4cm}} \hfill 日期：\underline{\hspace{3cm}}
        % 4. 学位论文原创性声明
% ============================================================
% frontmatter/supervisor_commitment.tex —— 研究生学位论文导师承诺书
% 通用模板，请以学校下发的最新版本替换正文。
% ============================================================
\chapter*{研究生学位论文导师承诺书}
\addcontentsline{toc}{chapter}{研究生学位论文导师承诺书}

本人作为该学位论文的指导教师，郑重承诺：
所指导的学位论文是在本人的指导下，由研究生独立完成的原创性研究工作；
论文内容真实、数据可靠，不存在抄袭、剽窃、伪造等学术不端行为；
论文中的研究成果未一稿多投或重复发表；
本人对论文的学术规范和质量负责。

\vspace{1.5cm}
\noindent 导师签名：\underline{\hspace{4cm}} \hfill 日期：\underline{\hspace{3cm}}
 % 5. 研究生学位论文导师承诺书
% ============================================================
% frontmatter/authorization.tex —— 学位论文使用授权书
% 文字已按《中国地质大学（武汉）研究生学位论文写作规范》（2025.5）
%   附件 5 官方式样核校（与官方一致）。
% ============================================================
\chapter*{学位论文使用授权书}
\addcontentsline{toc}{chapter}{学位论文使用授权书}

本人授权中国地质大学（武汉）可采用影印、缩印、数字化或其它复制手段保存本学位论文；学校可向国家有关部门或机构送交本学位论文的电子版全文，编入有关数据库进行检索、下载及文献传递服务；同意在校园网内提供全文浏览和下载服务。涉密论文解密后适用于本授权书。

\vspace{1.5cm}
\noindent 学位论文作者签名：\underline{\hspace{4cm}} \hfill 指导教师签名：\underline{\hspace{3cm}} \\
\noindent 日期：\underline{\hspace{3cm}} \hfill 日期：\underline{\hspace{3cm}}
      % 6. 学位论文使用授权书
% ============================================================
% frontmatter/author_bio.tex —— 作者简介（附件 6）
% 版式严格按 2025.5 规范附件 6 官方式样：
%   页标题"作者简历"；一、基本情况（姓名/性别/民族/出生年月/籍贯单行，
%   下接学历行）；二、学术论文；三、获奖、专利情况；四、研究项目。
% 注意：官方附件 6 无学号/学位类别/学科专业/研究方向/导师栏目
%   （这些信息已在封面与题名页体现），请勿再加回表格。
% ============================================================
\chapter*{作者简历}
\addcontentsline{toc}{chapter}{作者简历}

\vspace{0.5em}
\noindent 一、基本情况

\vspace{0.5em}
\noindent 姓名：\thesisAuthor\hspace{1em}性别：\underline{\hspace{1.1cm}}\hspace{1em}民族：\underline{\hspace{1.1cm}}\hspace{1em}出生年月：\underline{\hspace{1.6cm}}\hspace{1em}籍贯：\underline{\hspace{1.8cm}}

\vspace{0.5em}
\noindent 20XX 年 9 月—20XX 年 6 月\quad \thesisOrg\ \thesisMajor\ 攻读硕士学位，师从\thesisAdvisor。

\vspace{1.2em}
\noindent 二、学术论文

\vspace{0.5em}
\noindent 1. 作者. 题名[J]. 刊名，年，卷（期）：起-止页码.

\vspace{1.2em}
\noindent 三、获奖、专利情况

\vspace{0.5em}
\noindent 1. 获奖 / 专利名称（如有），等级，排名.

\vspace{1.2em}
\noindent 四、研究项目

\vspace{0.5em}
\noindent 1. 项目名称（基金 / 课题），项目时间，项目编号，角色.
         % 7. 作者简介
% ============================================================
% frontmatter/defense_committee.tex —— 学位论文答辩委员会名单（附件 7）
% 版式严格按 2025.5 规范附件 7 官方式样（Word 实测坐标，版心 [85.04, 510.24]）：
%   标题"学位论文答辩委员会名单"（黑体居中）；
%   无边框对齐四列：职务  姓名  职称  单位；
%   主席 1 行 + 委员 4 行 + 秘书 1 行 = 6 行（职务只在组内首行出现）；
%   式样底部无"答辩日期"栏，请勿添加。
%
% 【P1-4】旧版只有 4 人（主席 1 + 委员 2 + 秘书 1）。附件 7 式样与师兄 p7 均为
%   主席 1 + 委员 4 + 秘书 1 = 6 行，此处补齐。
% 【P2-8】列位按附件 7 式样：职务 r0=0.091 / 姓名 r0=0.298 / 职称 r0=0.468 / 单位 r0=0.694。
% 【P2-7】行距：式样 31.2pt、师兄 p7 实测 32.5pt —— 取 32pt。
% 【用户 2026-09-29 更正】单位列必须写全称「中国地质大学（武汉）」（师兄 p7 即为全称）。
%   全称比简称多「（武汉）」约 32pt，原 0.306 宽单位盒放不下会越出版心右缘。
%   为在不改字号(三号,规范)的前提下仍收进版心：把「职称」盒由 0.226 收到 0.156、
%   「单位」盒由 0.306 增到 0.376（单位起点随之由 0.694 调到 0.624，仍在版心内），
%   全称实测落在 x≈265–427pt，右端 427 < 版心右 510，整齐收边。
%
% ⚠️⚠️⚠️ 三个已踩过的坑（务必保留本注释，别再回退）：
%  1) 列宽**不能用 `p{...}`**：XeLaTeX 下 `p` 列会变成"定宽＋内容水平居中＋垂直居中"，
%     6 行内容全叠印在同一坐标（pdftoppm 实测 6 行文字全落在 y=231.2pt）。
%  2) `\makebox[宽][l]{...}` 产生的盒子**高度为 0**，放进 tabular 里该行
%     **不产生垂直步进**，同样导致 6 行叠印（PDF 内容流实测只有 Td 水平位移）。
%  3) `tabular` 里写 `\\[32pt]` 只加"行间额外距离"，且被 \arraystretch/\baselineskip
%     的交互扯乱，仍会叠行。
% ⇒ 最终方案：**彻底不用 tabular**，改成 6 行独立的 `\makebox[\textwidth][l]{...}`
%   用 `\\[32pt]` 手动分行（每行一个段落行，步进可靠），列位置用固定宽度
%   `\makebox[..][l]{}` 逐列锚定。实测 6 行分别落在 y=224/256/288/320/352/384pt，
%   步进 32pt，列位 0.091/0.298/0.468/0.694 全部命中。
% ============================================================
\chapter*{学位论文答辩委员会名单}
\addcontentsline{toc}{chapter}{学位论文答辩委员会名单}

\vspace{1.5cm}
\begin{center}
  \zihao{3}\songti
  % 注意：\makebox 高度为 0，若直接换行排，行与行会贴在一起（步进=0）。
  % 因此每行末尾都用 \rule{0pt}{32pt} 显式撑出行高，步进即 32pt。
  \lineskip=0pt\lineskiplimit=0pt\baselineskip=32pt\relax
  \noindent\makebox[\textwidth][l]{%
    \hspace*{0.091\textwidth}%
    \makebox[0.207\textwidth][l]{主席}%
    \makebox[0.170\textwidth][l]{×××}%
    \makebox[0.140\textwidth][l]{教授}%
    \makebox[0.392\textwidth][l]{中国地质大学（武汉）}\rule{0pt}{32pt}}\par
  \noindent\makebox[\textwidth][l]{%
    \hspace*{0.091\textwidth}%
    \makebox[0.207\textwidth][l]{委员}%
    \makebox[0.170\textwidth][l]{×××}%
    \makebox[0.140\textwidth][l]{副教授}%
    \makebox[0.392\textwidth][l]{中国地质大学（武汉）}\rule{0pt}{32pt}}\par
  \noindent\makebox[\textwidth][l]{%
    \hspace*{0.091\textwidth}%
    \makebox[0.207\textwidth][l]{}%
    \makebox[0.170\textwidth][l]{×××}%
    \makebox[0.140\textwidth][l]{副教授}%
    \makebox[0.392\textwidth][l]{中国地质大学（武汉）}\rule{0pt}{32pt}}\par
  \noindent\makebox[\textwidth][l]{%
    \hspace*{0.091\textwidth}%
    \makebox[0.207\textwidth][l]{}%
    \makebox[0.170\textwidth][l]{×××}%
    \makebox[0.140\textwidth][l]{副教授}%
    \makebox[0.392\textwidth][l]{中国地质大学（武汉）}\rule{0pt}{32pt}}\par
  \noindent\makebox[\textwidth][l]{%
    \hspace*{0.091\textwidth}%
    \makebox[0.207\textwidth][l]{}%
    \makebox[0.170\textwidth][l]{×××}%
    \makebox[0.140\textwidth][l]{副教授}%
    \makebox[0.392\textwidth][l]{中国地质大学（武汉）}\rule{0pt}{32pt}}\par
  \noindent\makebox[\textwidth][l]{%
    \hspace*{0.091\textwidth}%
    \makebox[0.207\textwidth][l]{秘书}%
    \makebox[0.170\textwidth][l]{×××}%
    \makebox[0.140\textwidth][l]{讲师}%
    \makebox[0.392\textwidth][l]{中国地质大学（武汉）}\rule{0pt}{32pt}}\par
\end{center}
  % 8. 学位论文答辩委员会名单
% ============================================================
% frontmatter/abstract.tex —— 中英文摘要
% 标题/关键词格式见 setup/format.tex 第 6 节；
% 正文为宋体小四(中文) / Times New Roman 小四(英文)，固定 20 磅行距。
% ============================================================

% ---------------- 中文摘要 ----------------
\clearpage
\cnabstracttitle
\addcontentsline{toc}{chapter}{摘要}    % 如需目录中显示摘要，保留此行；否则删除

% —— 中文摘要正文（示例，替换为真实内容）——
针对广域输电网常态化动态运维巡检场景，研究多机巢起降条件下的多无人机动态协同任务规划问题。
本文构建"动态任务选择 + 选择性多无人机协同任务分配（MUAS）"的两阶段分层求解框架，
并嵌入滚动时域运行机制，兼顾求解实时性与资源利用效率。

\vspace{12pt}
\noindent\textbf{关键词：}多无人机；任务分配；航迹规划；滚动时域；差分进化

% ---------------- 英文摘要 ----------------
\clearpage
% 【P2-9】去掉英文摘要页眉。依据：师兄论文 p9/p10（摘要/Abstract）实测 y>740 区域
%   完全为空 —— 既无页眉也无页脚；规范 2.4「页眉从第一章开始」。
%   规范 3.1 的"Abstract 部分用 Times New Roman"指正文用 TNR 字体，非要求加页眉。
%   （\fancypagestyle{abstracten} 定义保留在 setup/format.tex，需要时可再启用。）
\enabstracttitle
\addcontentsline{toc}{chapter}{Abstract}  % 如需目录中显示英文摘要，保留此行；否则删除

% —— English abstract（与中文摘要对应）——
% 【P2-10】师兄 p10 只有 Abstract + 正文，正文前不再重复排一遍英文题名（原与页眉重复）。
Aiming at the routine dynamic operation and maintenance inspection of wide-area
power transmission networks, this thesis studies the dynamic cooperative mission
planning of multiple unmanned aerial vehicles (UAVs) under multi-nest launch and
recovery conditions. A two-stage hierarchical framework combining dynamic task
selection and selective multi-UAV cooperative task allocation (MUAS) is proposed,
embedded in a rolling-horizon mechanism to balance solution timeliness and resource
utilization efficiency.

\vspace{12pt}
\noindent\textbf{Key Words:} multi-UAV; task allocation; path planning; rolling horizon; differential evolution

% 英文摘要页沿用前置默认 plain 样式（页脚居中罗马页码、无页眉），无需再显式切换。
           % 9. 中文摘要 / 10. Abstract
\tableofcontents                       % 11. 目录
% 12. 图和表清单（规范：如有）。当前正文尚无图/表，若编译出空目录页则保持注释；
% 插图/表后取消下面两行注释即可启用（并确认目录中不出现"图清单/表清单"以外的异常）。
% \listoffigures
% \listoftables

% ============================================================
% 正文部分（每章一个独立文件，便于维护）
% ============================================================
\mainmatter
\pagestyle{fancy}              % 正文起：启用页眉（中国地质大学格式）+ 外侧阿拉伯页码
% ============================================================
% chapters/ch01_intro.tex —— 第 1 章 绪论
% ============================================================
\chapter{绪论}

\section{研究背景与意义}
% 广域输电网巡检需求、无人机规模化应用、智能化调度必要性

\section{国内外研究现状}
\subsection{多无人机任务分配研究}
\subsection{航迹规划研究}
\subsection{动态协同调度研究}

\section{现有研究不足}
% 静态假设、任务集合固定、缺乏多机巢与滚动机制等

\section{本文主要研究内容}
% 任务分配、航迹规划、系统集成三条主线

\section{本文组织结构}
% 各章安排

% 提示：图表、公式、引用示例
% \begin{figure}[htbp]
%   \centering
%   \includegraphics[width=0.6\textwidth]{figures/xxx.pdf}
%   \caption{示意图}\label{fig:xxx}
% \end{figure}

% ============================================================
% chapters/ch02_related.tex —— 第 2 章 相关理论与技术
% ============================================================
\chapter{相关理论与技术}

\section{多无人机协同调度问题建模}
\subsection{任务分配问题（MTSP / VRP 视角）}
\subsection{航迹规划问题（三维代价场）}

\section{滚动时域优化方法}
% Rolling Horizon / Model Predictive Control 基本原理

\section{进化算法基础}
\subsection{差分进化算法}
\subsection{离散映射与编码机制}

\section{三维地理环境与可达性建模}
% 地形、禁飞区、三维执行代价矩阵

\section{本章小结}

% ============================================================
% chapters/ch03_allocation.tex —— 第 3 章 动态滚动的多机巢多无人机协同任务分配方法
% 小节框架已按既有方法论（task_allocation/Chapter3_README.md）拆好，
% 请逐节补充正文、公式与图示。
% ============================================================
\chapter{动态滚动的多机巢多无人机协同任务分配方法}

\section{引言}
% 本章面向广域输电网动态运维巡检，研究多机巢起降条件下的多无人机动态协同任务规划问题。

\section{问题特征与核心挑战}
\subsection{任务动态到达与任务持续滞留}
\subsection{无人机可用状态动态变化}
\subsection{多机巢起降与终点机巢选择}
\subsection{任务数量与当前执行能力不匹配}
\subsection{任务价值与空间执行代价并存}

\section{总体求解框架}
\subsection{第一阶段：动态任务选择}
\subsection{第二阶段：选择性 MUAS}

\section{第一阶段：动态任务选择机制}
\subsection{任务选择的基本思想}
\subsection{基础任务优先级}
\subsection{候选任务预筛选}
\subsection{边际收益增量选择}
\subsection{快速可行性检查}

\section{第二阶段：选择性多机巢 MUAS 优化}
\subsection{问题定义}
\subsection{决策变量}
\subsection{多目标优化模型}

\section{改进离散映射差分进化算法}
\subsection{面向任务序列的统一编码}
\subsection{多机巢终点选择}
\subsection{选择性任务分配}
\subsection{约束处理}
\subsection{三维地理代价嵌入}

\section{动态滚动与闭环运行}

\section{本章研究内容与贡献定位}
\subsection{构建动态任务选择机制}
\subsection{构建选择性多机巢 MUAS 模型及改进求解方法}
\subsection{形成动态任务选择与精细 MUAS 相结合的分层求解框架}

\section{本章总体技术路线}

\section{本章小结}

% 提示：算法示例
% \begin{algorithm}[htbp]
%   \caption{选择性多机巢 MUAS 求解}\label{alg:muas}
%   \begin{algorithmic}[1]
%     \Require 候选任务集 $\mathcal{T}_t^{\mathrm{sel}}$，可用无人机集 $U_t^{\mathrm{avail}}$
%     \Ensure 分配方案 $\mathcal{T}_t^{\mathrm{exec}}, U_t^{\mathrm{exec}}$
%     \For{每个滚动周期 $t$}
%       \State 动态任务选择 $\rightarrow \mathcal{T}_t^{\mathrm{sel}}$
%       \State 改进离散映射差分进化求解 MUAS
%     \EndFor
%     \Return 最优分配与航次
%   \end{algorithmic}
% \end{algorithm}

% ============================================================
% chapters/ch04_planning.tex —— 第 4 章 顾及风险时空异质性的航迹规划方法
% ============================================================
\chapter{顾及风险时空异质性的航迹规划方法}

\section{引言}

\section{风险时空异质性建模}
\subsection{威胁源时空分布}
\subsection{风险代价场构建}

\section{三维航迹表示与约束}
\subsection{动力学约束}
\subsection{禁飞区与可达性约束}

\section{航迹优化方法}
\subsection{代价函数设计}
\subsection{搜索 / 优化策略}

\section{与任务分配的耦合}

\section{仿真实验与结果分析}

\section{本章小结}

% ============================================================
% chapters/ch05_system.tex —— 第 5 章 智能调度系统实现与应用验证
% 可结合仓库 task_allocation/ 下的实验框架（exp01_smoke 等）撰写。
% ============================================================
\chapter{智能调度系统实现与应用验证}

\section{引言}

\section{系统总体架构}
% 数据层 / 算法层 / 实验编排层（参考 task_allocation/experiments/README.md 约定）

\section{关键模块实现}
\subsection{动态任务池管理}
\subsection{任务选择模块}
\subsection{MUAS 求解模块}
\subsection{滚动与事件触发模块}

\section{实验设计与参数设置}
% 对应 task_allocation/experiments/exp0X_*/config.yaml

\section{应用验证与结果分析}

\section{本章小结}

% ============================================================
% chapters/ch06_conclusion.tex —— 第 6 章 总结与展望
% ============================================================
\chapter{总结与展望}

\section{全文工作总结}
% 逐章回顾主要工作与结论

\section{主要创新点}
% 1) 动态任务选择机制；2) 选择性多机巢 MUAS 及改进算法；3) 分层滚动框架

\section{研究不足与未来展望}
\subsection{算法规模与实时性}
\subsection{不确定性建模}
\subsection{系统落地与工程化}

\section{本章小结}


% ============================================================
% 后置部分（按学位论文组成顺序：致谢→参考文献→附录→攻读学位期间成果）
% ============================================================
\backmatter
% ============================================================
% backmatter/acknowledgements.tex —— 致谢
% ============================================================
\chapter*{致谢}
\addcontentsline{toc}{chapter}{致谢}

% 在此感谢导师、家人、同学与项目资助等。
    % 14. 致谢
\printbibliography[heading=bibintoc, title={参考文献}]  % 15. 参考文献
% ============================================================
% backmatter/appendices.tex —— 附录
% 每个附录用 \appendix + \chapter{} 或 \section{} 组织
% ============================================================
\appendix
\chapter{证明与推导}
% 定理证明、复杂公式推导等

\chapter{补充实验数据}
% 主文中未展开的图表、参数表
          % 16. 附录
% ============================================================
% backmatter/publications.tex —— 攻读学位期间科研成果（可选）
% ============================================================
\chapter*{攻读学位期间科研成果}
\addcontentsline{toc}{chapter}{攻读学位期间科研成果}

\section*{发表的学术论文}
% [1] 作者. 题名[J]. 刊名, 年份, 卷(期): 页码.

\section*{参与的科研项目}
% 项目名称、编号、角色
        % 攻读学位期间科研成果（可选）

\end{document}
